%\title{Bechelor Thesis Your namet}
\RequirePackage[l2tabu]{nag}		% Warns for incorrect (obsolete) LaTeX usage
%
%
% File: John Doe - Your Matriculation number
% Author: John Doe
% Description: Your theis topic
% 
% Memoir is a flexible class for typesetting poetry, fiction, 
% non-fiction and mathematical works as books, reports, articles or
% manuscripts. CTAN repository is found at:
% http://www.ctan.org/tex-archive/macros/latex/contrib/memoir/
%
%
% HSG guidelines for thesis presentation were found at:
% place a link that is missing here 
%
%
% Memoir class loads useful packages by default (see manual).
\documentclass[a4paper,11pt,leqno,openbib,11pt,openany,oneside]{memoir} %add 'draft' to turn draft option on (see below)
%
\linespread{1.2}
% Adding metadata:
\usepackage{datetime}
\usepackage{ifpdf}
\ifpdf
\pdfinfo{
   /Author (John Doe)
   /Title (Bachelor Thesis)
   /Keywords (Stock Analysis; Time series; Whatever)
   /CreationDate (D:\pdfdate)
}
\fi
% When draft option is on. 
\ifdraftdoc 
	\usepackage{draftwatermark}				%Sets watermarks up.
	\SetWatermarkScale{0.3}
	\SetWatermarkText{\bf Draft: \today}
\fi
%
% Declare figure/table as a subfloat.
\newsubfloat{figure}
\newsubfloat{table}
% Better page layout for A4 paper, see memoir manual.
\settrimmedsize{297mm}{210mm}{*}
\setlength{\trimtop}{0pt} 
\setlength{\trimedge}{\stockwidth} 
\addtolength{\trimedge}{-\paperwidth} 
\settypeblocksize{648pt}{448.13pt}{*} 
\setulmargins{3.0cm}{*}{*} 
\setlrmargins{*}{*}{1.5} 
\setmarginnotes{17pt}{51pt}{\onelineskip} 
\setheadfoot{\onelineskip}{2\onelineskip} 
\setheaderspaces{*}{11pt}{*} 
\checkandfixthelayout
%
\frenchspacing
% Font with math support: New Century Schoolbook
\usepackage{fouriernc}
\usepackage[T1]{fontenc}
% package for tables 
\usepackage{tabularx}
% package for landscape
\usepackage{pdflscape}
\usepackage{fancyhdr} 

%\fancypagestyle{mylandscape}{
%\fancyhf{} %Clears the header/footer
%\fancyfoot{% Footer
%\makebox[\textwidth][r]{% Right
%  \rlap{\hspace{.75cm}% Push out of margin by \footskip
 %   \smash{% Remove vertical height
  %    \raisebox{4.87in}{% Raise vertically
  %      \rotatebox{90}{\thepage}}}}}}% Rotate counter-clockwise
%\renewcommand{\headrulewidth}{0pt}% No header rule
%\renewcommand{\footrulewidth}{0pt}% No footer rule
%}

% Current guidelines:
%
% Text should be in double or 1.5 line spacing, and font size should be
% chosen to ensure clarity and legibility for the main text and for any
% quotations and footnotes. Margins should allow for eventual hard binding.
%
% Note: This is automatically set by memoir class. Nevertheless \OnehalfSpacing 
% enables double spacing but leaves single spaced for captions for instance. \OnehalfSpacing 
%
% Sets numbering division level
\setsecnumdepth{subsection} 
\maxsecnumdepth{subsubsubsection}
%
% Chapter style (taken and slightly modified from Lars Madsen Memoir Chapter 
% Styles document
\usepackage{calc,soul,fourier}
\makeatletter 
\newlength\dlf@normtxtw 
\setlength\dlf@normtxtw{\textwidth} 
\newsavebox{\feline@chapter} 
\newcommand\feline@chapter@marker[1][4cm]{%
	\sbox\feline@chapter{% 
		\resizebox{!}{#1}{\fboxsep=1pt%
			\colorbox{gray}{\color{white}\thechapter}% 
		}}%
		\rotatebox{90}{% 
			\resizebox{%
				\heightof{\usebox{\feline@chapter}}+\depthof{\usebox{\feline@chapter}}}% 
			{!}{\scshape\so\@chapapp}}\quad%
		\raisebox{\depthof{\usebox{\feline@chapter}}}{\usebox{\feline@chapter}}%
} 
\newcommand\feline@chm[1][4cm]{%
	\sbox\feline@chapter{\feline@chapter@marker[#1]}% 
	\makebox[0pt][c]{% aka \rlap
		\makebox[1cm][r]{\usebox\feline@chapter}%
	}}
\makechapterstyle{daleifmodif}{
	\renewcommand\chapnamefont{\normalfont\Large\scshape\raggedleft\so} 
	\renewcommand\chaptitlefont{\normalfont\Large\bfseries\scshape} 
	\renewcommand\chapternamenum{} \renewcommand\printchaptername{} 
	\renewcommand\printchapternum{\null\hfill\feline@chm[2.5cm]\par} 
	\renewcommand\afterchapternum{\par\vskip\midchapskip} 
	\renewcommand\printchaptertitle[1]{\color{gray}\chaptitlefont\raggedleft ##1\par}
} 
\makeatother 
\chapterstyle{daleifmodif}
%
% Current guidelines:
%
% The pages should be numbered consecutively at the bottom centre of the
% page.
\makepagestyle{myvf} 
\makeoddfoot{myvf}{}{\thepage}{} 
\makeevenfoot{myvf}{}{\thepage}{} 
\makeheadrule{myvf}{\textwidth}{\normalrulethickness} 
\makeevenhead{myvf}{\small\textsc{\leftmark}}{}{} 
\makeoddhead{myvf}{}{}{\small\textsc{\rightmark}}
\pagestyle{myvf}
%
% Oscar's command (it works):
% Fills blank pages until next odd-numbered page. Used to emulate single-sided
% frontmatter. This will work for title, abstract and declaration. Though the
% contents sections will each start on an odd-numbered page they will
% spill over onto the even-numbered pages if extending beyond one page
% (hopefully, this is ok).
\newcommand{\clearemptydoublepage}{\newpage{\thispagestyle{empty}\cleardoublepage}}
%
%
% Creates indexes for Table of Contents, List of Figures, List of Tables and Index
\makeindex
% \printglossaries below creates a list of abbreviations. \gls and related
% commands are then used throughout the text, so that latex can automatically
% keep track of which abbreviations have already been defined in the text.
%
% The import command enables each chapter tex file to use relative paths when
% accessing supplementary files. For example, to include
% chapters/brewing/images/figure1.png from chapters/brewing/brewing.tex we can
% use
% \includegraphics{images/figure1}
% instead of
% \includegraphics{chapters/brewing/images/figure1}
\usepackage{import}

% Add other packages needed for chapters here. For example:
\usepackage{lipsum}					%Needed to create dummy text
%\usepackage{subfig}                 %Two pictures side-by-side
\usepackage{subcaption} 

\usepackage{siunitx}       %Center table values
\newcolumntype{L}{>{\phantom{$\mathbin{-}$}$}l<{$}}
\newcommand{\mi}{\mathllap{\mathbin{-}}}

\usepackage{amsfonts} 					%Calls Amer. Math. Soc. (AMS) fonts
\usepackage[centertags,leqno]{amsmath}  %Writes maths centred down
\makeatletter
\newcommand{\leqnomode}{\tagsleft@true}
\newcommand{\reqnomode}{\tagsleft@false}
\makeatother
\usepackage{amssymb,amsthm}
\usepackage{thmtools}
\declaretheoremstyle[
spaceabove=6pt, spacebelow=6pt,
headfont=\normalfont\bfseries,
notefont=\mdseries, notebraces={(}{)},
bodyfont=\normalfont,
postheadspace=0.6em,
headpunct=:
]{mystyle}
\declaretheorem[style=mystyle, name=Hypothesis, preheadhook={\renewcommand{\thehyp}{H\textsubscript{\arabic{hyp}}}}]{hyp}


%%%%%%%%%%%%%%%%%%%%%%%%%%%%%%%%%%%%%%%%%%%%%%%%%%%
\usepackage{acro}

\DeclareAcronym{ate}{
  short=ATE,
  long=Average Treatment Effect,
}
\DeclareAcronym{att}{
  short=ATT,
  long=Average Treatment Effect on the Treated,
}
\DeclareAcronym{hte}{
  short=HTE,
  long=Heterogeneous Treatment Effect,
}
\DeclareAcronym{cate}{
  short=CATE,
  long=Conditional Average Treatment Effect,
}
\DeclareAcronym{SEK}{
  short=SEK,
  long=Swedish Krona,
}
\DeclareAcronym{GBP}{
  short=GBP,
  long=Green Bond Principles,
}
\DeclareAcronym{CBI}{
  short=CBI,
  long=Climate Bonds Initiative,
}
\DeclareAcronym{OLS}{
  short=OLS,
  long=Ordinary Least Squares,
}
\DeclareAcronym{CF}{
  short=CF,
  long=Causal Forest,
}
\DeclareAcronym{PSM}{
  short=PSM,
  long=Propensity Score Matching,
}
\DeclareAcronym{CEM}{
  short=CEM,
  long=Coarsened Exact Matching,
}
\DeclareAcronym{CIA}{
  short=CIA,
  long=Conditional Independence Assumption,
}
\DeclareAcronym{AIPW}{
  short=AIPW,
  long=Augmented Inverse Probability Weighting,
}
\DeclareAcronym{SUTVA}{
  short=SUTVA,
  long=Stable Unit Treatment Value Assumption,
}
\DeclareAcronym{i.i.d}{
  short=i.i.d,
  long=identically and independently distributed,
}
\DeclareAcronym{DGP}{
  short=DGP,
  long=Data Generating Process,
}

%%%%%%%%%%%%%%%%%%%%%%%%%%%%%%%%%%%%%%%%%%%%%%%%%%%
%%%%%%%%%%%%%%%%%%%%%%%%%%%%%%%%%%%%%%%%%%%%%%%%%%%
\usepackage{array}
%%%%%%%%%%%%%%%%%%%%%%%%%%%%%%%%%%%%%%%%%%%%%%%%%%%
%%%%%%%%%%%%%%%%%%%%%%%%%%%%%%%%%%%%%%%%%%%%%%%%%%%
%%%%%%%%%%%%%%%%%%%%%%%%%%%%%%%%%%%%%%%%%%%%%%%%%%%
\usepackage[acronym]{glossaries}
\makeglossaries
\usepackage{stmaryrd}					%New AMS symbols
\usepackage{amssymb}					%Calls AMS symbols
\usepackage{amsthm}					%Calls AMS theorem environment
\usepackage{newlfont}					%Helpful package for fonts and symbols
\usepackage{layouts}					%Layout diagrams
\usepackage{graphicx}					%Calls figure environment
\usepackage{longtable,rotating}			%Long tab environments including rotation. 
\usepackage[applemac]{inputenc}			%Needed to encode non-english characters 
									%directly for mac
\usepackage{colortbl}					%Makes coloured tables
\usepackage[x11names]{xcolor}
\usepackage{threeparttable}             %Footnote after table
\usepackage{wasysym}					%More math symbols
\usepackage{mathrsfs}					%Even more math symbols
\usepackage{float}						%Helps to place figures, tables, etc. 
\usepackage{placeins}
\usepackage{verbatim}					%Permits pre-formated text insertion
\usepackage{upgreek}					%Calls other kind of greek alphabet
\usepackage{latexsym}					%Extra symbols
\usepackage[authoryear,square,semicolon,sort&compress]{natbib}

\definecolor{citecolour}{rgb}{0.0,0.5,0.0}
\RequirePackage[pdfusetitle]{hyperref}
\hypersetup{
     bookmarksnumbered=true,
     colorlinks=true,
     linkcolor=citecolour,
     citecolor=citecolour,
     filecolor=black,
     urlcolor=black
}%Calls bibliography commands 

\usepackage{listings}
\usepackage{rotating}
\usepackage{enumitem}
\usepackage{url}						%Supports url commands
\usepackage{csquotes}
\MakeOuterQuote{"}
\usepackage{etex}					 	%eTeXÕs extended support for counters
\usepackage{fixltx2e}					%Eliminates some in felicities of the 
									%original LaTeX kernel
\usepackage[spanish,english]{babel}		%For languages characters and hyphenation
\usepackage{color}                    				%Creates coloured text and background
\usepackage[colorlinks=true,
		     allcolors=black]{hyperref}              %Creates hyperlinks in cross references
\usepackage{cleveref}
\crefname{hyp}{hypothesis}{hypotheses}
\Crefname{hyp}{Hypothesis}{Hypotheses}
\usepackage{memhfixc}					%Must be used on memoir document 
									%class after hyperref
\usepackage{enumerate}					%For enumeration counter
\usepackage{footnote}					%For footnotes
\usepackage{microtype}					%Makes pdf look better.
\usepackage{rotfloat}					%For rotating and float environments as tables, 
									%figures, etc. 
\usepackage{alltt}						%LaTeX commands are not disabled in 
									%verbatim-like environment
\usepackage[version=0.96]{pgf}			%PGF/TikZ is a tandem of languages for producing vector graphics from a 
\usepackage{graphics}
\usepackage{adjustbox}
\usepackage{tikz}						%geometric/algebraic description.
\usetikzlibrary{arrows,shapes,snakes,
		       automata,backgrounds,
		       petri,topaths}				%To use diverse features from tikz		
%							
%Reduce widows  (the last line of a paragraph at the start of a page) and orphans 
% (the first line of paragraph at the end of a page)
\widowpenalty=1000
\clubpenalty=1000
%
% New command definitions for my thesis
%
\newcommand{\keywords}[1]{\par\noindent{\small{\bf Keywords:} #1}} %Defines keywords small section
\newcommand{\parcial}[2]{\frac{\partial#1}{\partial#2}}                             %Defines a partial operator
\newcommand{\vectorr}[1]{\mathbf{#1}}                                                        %Defines a bold vector
\newcommand{\vecol}[2]{\left(                                                                         %Defines a column vector
	\begin{array}{c} 
		\displaystyle#1 \\
		\displaystyle#2
	\end{array}\right)}
\newcommand{\mados}[4]{\left(                                                                       %Defines a 2x2 matrix
	\begin{array}{cc}
		\displaystyle#1 &\displaystyle #2 \\
		\displaystyle#3 & \displaystyle#4
	\end{array}\right)}
\newcommand{\pgftextcircled}[1]{                                                                    %Defines encircled text
    \setbox0=\hbox{#1}%
    \dimen0\wd0%
    \divide\dimen0 by 2%
    \begin{tikzpicture}[baseline=(a.base)]%
        \useasboundingbox (-\the\dimen0,0pt) rectangle (\the\dimen0,1pt);
        \node[circle,draw,outer sep=0pt,inner sep=0.1ex] (a) {#1};
    \end{tikzpicture}
}
\newcommand{\range}[1]{\textnormal{range }#1}                                             %Defines range operator
\newcommand{\innerp}[2]{\left\langle#1,#2\right\rangle}                                 %Defines inner product
\newcommand{\prom}[1]{\left\langle#1\right\rangle}                                         %Defines average operator
\newcommand{\tra}[1]{\textnormal{tra} \: #1}                                                       %Defines trace operator
\newcommand{\sign}[1]{\textnormal{sign\,}#1}                                                   %Defines sign operator
\newcommand{\sech}[1]{\textnormal{sech} #1}                                                  %Defines sech
\newcommand{\diag}[1]{\textnormal{diag} #1}                                                    %Defines diag operator
\newcommand{\arcsech}[1]{\textnormal{arcsech} #1}                                       %Defines arcsech
\newcommand{\arctanh}[1]{\textnormal{arctanh} #1}                                         %Defines arctanh
%Change tombstone symbol
\newcommand{\blackged}{\hfill$\blacksquare$}
\newcommand{\whiteged}{\hfill$\square$}
\newcounter{proofcount}
\renewenvironment{proof}[1][\proofname.]{\par
 \ifnum \theproofcount>0 \pushQED{\whiteged} \else \pushQED{\blackged} \fi%
 \refstepcounter{proofcount}
 \normalfont 
 \trivlist
 \item[\hskip\labelsep
       \itshape
   {\bf\em #1}]\ignorespaces
}{%
 \addtocounter{proofcount}{-1}
 \popQED\endtrivlist
}
%
%
% New definition of square root:
% it renames \sqrt as \oldsqrt
\let\oldsqrt\sqrt
% it defines the new \sqrt in terms of the old one
\def\sqrt{\mathpalette\DHLhksqrt}
\def\DHLhksqrt#1#2{%
\setbox0=\hbox{$#1\oldsqrt{#2\,}$}\dimen0=\ht0
\advance\dimen0-0.2\ht0
\setbox2=\hbox{\vrule height\ht0 depth -\dimen0}%
{\box0\lower0.4pt\box2}}
%
% My caption style
\newcommand{\mycaption}[2][\@empty]{
	\captionnamefont{\scshape} 
	\changecaptionwidth
	\captionwidth{0.9\linewidth}
	\captiondelim{.\:} 
	\indentcaption{0.75cm}
	\captionstyle[\centering]{}
	\setlength{\belowcaptionskip}{10pt}
	\ifx \@empty#1 \caption{#2}\else \caption[#1]{#2}
}
%
% My subcaption style
\newcommand{\mysubcaption}[2][\@empty]{
	\subcaptionsize{\small}
	\hangsubcaption
	\subcaptionlabelfont{\rmfamily}
	\sidecapstyle{\raggedright}
	\setlength{\belowcaptionskip}{10pt}
	\ifx \@empty#1 \subcaption{#2}\else \subcaption[#1]{#2}
}
%
%An initial of the very first character of the content
\usepackage{lettrine}
\newcommand{\initial}[1]{%
	\lettrine[lines=3,lhang=0.33,nindent=0em]{
		\color{citecolour}
     		{\textsc{#1}}}{}}
%
% Theorem styles used in my thesis
%
\theoremstyle{plain}
\newtheorem{theo}{Theorem}[chapter]
\theoremstyle{plain}
\newtheorem{prop}{Proposition}[chapter]
\theoremstyle{plain}
\theoremstyle{definition}
\newtheorem{dfn}{Definition}[chapter]
\theoremstyle{plain}
\newtheorem{lema}{Lemma}[chapter]
\theoremstyle{plain}
\newtheorem{cor}{Corollary}[chapter]
\theoremstyle{plain}
\newtheorem{resu}{Result}[chapter]
%
% Hyphenation for some words
%
\hyphenation{res-pec-tively}
\hyphenation{mono-ti-ca-lly}
\hyphenation{hypo-the-sis}
\hyphenation{para-me-ters}
\hyphenation{sol-va-bi-li-ty}
%


%
\begin{document}
% Current guidlines:
%
% Preliminary pages
% 
% The five preliminary pages must be the Title Page, Abstract, Dedication
% and Acknowledgements, Author's Declaration and Table of Contents.
% These should be single-sided.
% 
% Table of contents, list of tables and illustrative material
% 
% The table of contents must list, with page numbers, all chapters,
 % sections and subsections, the list of references, bibliography, list of
% abbreviations and appendices. The list of tables and illustrations
% should follow the table of contents, listing with page numbers the
% tables, photographs, diagrams, etc., in the order in which they appear
% in the text.
% 
\frontmatter
\pagenumbering{roman}
%
%
% File: Title.tex
% Author: Your Name
% Description: Contains the title page
%

\begin{titlingpage}
\begin{SingleSpace}
\calccentering{\unitlength} 
\begin{adjustwidth*}{\unitlength}{-\unitlength}
\vspace*{1mm}
\begin{center}
\rule[0.5ex]{\linewidth}{2pt}\vspace*{-\baselineskip}\vspace*{3.2pt}
\rule[0.5ex]{\linewidth}{1pt}\\[\baselineskip]
{\HUGE \color{citecolour} The Green Bonds Premium Puzzle \par}\\[4mm]
{\Large \textit{Evidence from Causal Machine Learning}}\\
\rule[0.5ex]{\linewidth}{1pt}\vspace*{-\baselineskip}\vspace{3.2pt}
\rule[0.5ex]{\linewidth}{2pt}\\
\vspace{6.5mm}
{\large \textsc{Master's Thesis}}\\
\vspace{6.5mm}
{\large By}\\
\vspace{6.5mm}
{\large\textsc{Arbian Halilaj}\\
ID: 16-609-828}\\
\vspace{6mm}
\includegraphics[scale=0.35]{chinchilab-template/logos/unisg_new.png}\\
\vspace{6mm}
{\large Master of Arts in Economics\\
\textsc{University of St.Gallen}}\\
\vspace{6mm}
{\large Under the supervision of \\
\textsc{Prof. Dr. Jana Mareckova}}\\
\vspace{6mm}
{\large Co-supervisor \\
\textsc{Prof. Dr. Lyudmila Grigoryeva}}\\
\vspace{9mm}
\begin{minipage}{10cm}
A thesis submitted to the University of St.Gallen in accordance with the requirements of the degree of \textsc{Master of Arts} in the Faculty of Economics.
\end{minipage}\\
\vspace{9mm}
{\large\textsc{November 2022}}
\vspace{12mm}
\end{center}
%\begin{flushright}
%{\small Word count: ten thousand and four}
%\end{flushright}
\end{adjustwidth*}
\end{SingleSpace}
\end{titlingpage}
%\clearemptydoublepage
%
%
% File: abstract.tex
% Author: Your Name
% Description: Contains the text for thesis abstract
%

\chapter*{Abstract}

\begin{center}
    {\HUGE \color{citecolour} The Green Bonds Premium Puzzle \par}
    \vspace{2.5mm}
    {\Large \textit{Evidence from Causal Machine Learning}}
    \vspace{6.5mm}
    
    {\centering\large\textsc{Arbian Halilaj}}
    \vspace{6.5mm}
\end{center}


%\begin{SingleSpace}
\initial{T}his paper measures the greenium effect by using a novel tool from the causal machine learning literature. By employing the causal forest algorithm, we aim to estimate the average treatment effect (ATE) as well as the conditional average treatment effect (CATE) of a bond labelled as green in the primary market. Our data universe comes from the Eikon Refinitiv database, from which we have derived three distinct data samples. We employ three different estimation strategies, including an ensemble method with propensity score matching, resulting in a total of 11 models. Overall, our results show that the greenium effect ranges from 18.6 to 103.2 basis points, with a central tendency of 43 basis points. This effect is statistically and economically significant. We further show that violating the i.i.d assumption has a small negative effect on the estimation accuracy of the average treatment effect of causal forests.
%\end{SingleSpace}

\vspace{14mm}

{\textbf{JEL Classification}} 
A13; G12; Q01; Q51; Q56.

{\large \textbf{\textsc{Keywords:}}} {\textsc{Sustainable Finance; Green Bonds; Greenium; Primary Market; Causal Inference; Causal Machine Learning; Causal Forest.}}


\clearpage
%\clearemptydoublepage
%
%%
% File: discretion.tex
% Author: Your name
% Description: Contains the Discretion page
%
% HSG guidelines:
%
% Declaration of Discretion 
%
\chapter*{Declaration of Discretion}
\begin{SingleSpace}
\begin{quote}
The Undersigned\\ 

hereby undertakes and warrants to treat any information obtained by the enterprise/administration concerned in strict confidence. In particular, he / she shall only permit people other than the referees to inspect his / her written work with the express consent of all the parties that have provided information. \\

The Undersigned hereby takes cognizance of the fact that the University of St.Gallen may check his / her work for any plagiarism with the help of a plagiarism software and that the undersigned shall have to notify the enterprise/administration surveyed accordingly.

\vspace{1.5cm}
\noindent
\hspace{-0.75cm}\textsc{SIGNED: .................................................... DATE: ..........................................}
\end{quote}
\end{SingleSpace}
\clearpage
%\clearemptydoublepage
%
%
% File: declaration.tex
% Author: Your Name
% Description: Contains the declaration page
%
% In accordance with HSG guidelines 2019:
%
%
\chapter*{Declaration of authorship}
\begin{SingleSpace}
\begin{quote}
I hereby declare

\begin{itemize}
    \item that I have written this thesis without any help from others and without the use of documents and aids other than those stated above;
    \item that I have mentioned all the sources used and that I have cited them correctly according to established academic citation rules;
    \item that I have acquired any immaterial rights to materials I may have used such as images or graphs, or that I have produced such materials myself;
    \item that the topic or parts of it are not already the object of any work or examination of another course unless this has been explicitly agreed on with the faculty member in advance and is referred to in the thesis;
    \item that I will not pass on copies of this work to third parties or publish them without the University's written consent if a direct connection can be established with the University of St.Gallen or its faculty members;
    \item that I am aware that my work can be electronically checked for plagiarism and that I hereby grant the University of St.Gallen copyright in accordance with the Examination Regulations in so far as this is required for administrative action;
    \item that I am aware that the University will prosecute any infringement of this declaration of authorship and, in particular, the employment of a ghostwriter, and that any such infringement may result in disciplinary and criminal consequences which may result in my expulsion from the University or my being stripped of my degree.
\end{itemize}


\vspace{1.5cm}
\noindent
\hspace{-0.75cm}\textsc{SIGNED: .................................................... DATE: ..........................................}\\

By submitting this academic term paper, I confirm through my conclusive action that I am
submitting the Declaration of Authorship, that I have read and understood it, and that it is
true.
\end{quote}
\end{SingleSpace}
\clearpage
%\clearemptydoublepage
%
\renewcommand{\contentsname}{Table of Contents}
\maxtocdepth{subsubsection}
\tableofcontents*
\addtocontents{toc}{\par\nobreak \mbox{}\hfill{\bf Page}\par\nobreak}
%\clearemptydoublepage
\newpage
%
\listoftables
\addtocontents{lot}{\par\nobreak\textbf{{\scshape Table} \hfill Page}\par\nobreak}
%\clearemptydoublepage
\newpage
%
\listoffigures
\addtocontents{lof}{\par\nobreak\textbf{{\scshape Figure} \hfill Page}\par\nobreak}
%\clearemptydoublepage
%
\printacronyms[name=Nomenclature]
%
\mainmatter
%
\import{chapters/chapter01/}{chap01.tex}
%\clearemptydoublepage
%
%
\import{chapters/chapter02/}{chap02.tex}
%\clearemptydoublepage
%
%
\import{chapters/chapter03/}{chap03.tex}
%\clearemptydoublepage
%
%
\import{chapters/chapter04/}{chap04.tex}
%\clearemptydoublepage
%
%
\import{chapters/chapter05/}{chap05.tex}
%\clearemptydoublepage

\import{chapters/chapter06/}{chap06.tex}
%\clearemptydoublepage
%
% Apparently the guidelines don't say anything about citations or
% bibliography styles so I guess we can use anything.
\backmatter
%\bibliographystyle{munich.bst} %german
%\bibliographystyle{agsm.bst}
%\bibliographystyle{regstud.bst}
%\bibliographystyle{aer.bst}

%\bibliographystyle{ecca.bst}

%\bibliographystyle{itaxpf.bst}
%\bibliographystyle{siam-letters.bst}
%\bibliographystyle{computational-complexity/cc.bst}
%\bibliographystyle{gost780.bst}
\bibliographystyle{biblio2.bst}
%\bibliographystyle{biochem.bst}


\refstepcounter{chapter}
\bibliography{Shkenca}
%\setlength{\bibhang}{1cm}

%\setcitestyle{authoryear}
%\printbibliography
%\clearemptydoublepage
%
%
% And the appendix goes here
\appendix
\import{chapters/appendices/}{app0A.tex}
%\clearemptydoublepage

\import{chapters/appendices/}{app0B.tex}
%\clearemptydoublepage

\import{chapters/appendices/}{app0C.tex}
%\clearemptydoublepage

\import{chapters/appendices/}{app0D.tex}
%\clearemptydoublepage

\import{chapters/appendices/}{app0E.tex}
%\clearemptydoublepage

\import{chapters/appendices/}{app0F.tex}
%\clearemptydoublepage

%
% Add index
%\printindex
%   
\end{document}